\documentclass[10pt,a4paper]{amsart}
\usepackage[margin=19mm]{geometry}
\usepackage{amsmath,amssymb,mathtools}
\usepackage[scr=rsfs,cal=euler]{mathalfa}
\usepackage{xfrac}
\usepackage[unicode,hidelinks,bookmarks=false]{hyperref}
\newcommand{\ie}{i.e.\ }
\newcommand{\cf}{cf.\ }
\newcommand{\resp}{respectively\ }
\newcommand{\Subsection}{\subsection}
\providecommand{\nobreakdash}{\nobreak\hbox{-}}
\setlength{\emergencystretch}{2em}
\multlinegap0pt
\usepackage{amssymb}
\usepackage[shortlabels]{enumitem}
\usepackage{xcolor}
\newtheorem{lemma}{Lemma}
\title{Amphichiral Knots: Odd Braid Index and Symmetry Classification in the Three-Braid Case}
\author{Hyungseok Jung}
\date{Source: arXiv:2609.12295v1 (CC BY 4.0)}
\begin{document}
\maketitle

\noindent\emph{Source and licence.} Amphichiral Knots: Odd Braid Index and Symmetry Classification in the Three-Braid Case. Version arXiv:2609.12295v1 (CC BY 4.0), distributed by arXiv under the Creative Commons Attribution 4.0 International licence, http://creativecommons.org/licenses/by/4.0/, as recorded in the arXiv licence metadata for this exact version and shown on the abstract page https://arxiv.org/abs/2609.12295v1. Full licence notice: \texttt{licenses/arxiv-2609.12295-CC-BY-4.0.txt}.\par\smallskip

\noindent\textit{Editorial scope.} Counted unit: the article's lemma lem:flype-identities (source0.tex lines 1368--1390). The article's complete original proof of it is delivered with it (source0.tex lines 1392--1434, one \texttt{proof} environment of 43 lines). No mathematics is written, repaired or re-derived: the statement and the proof are the article's own lines, in the article's own notation, and no retained line is substituted or re-flowed.  No further source-local context is needed: every cross-reference of every retained line resolves inside the two delivered spans.  Nothing was omitted from any retained span, so the package carries no omission record.  No retained line contains a citation, so the wrapper carries no bibliography and no bibliography entry.  No retained line contains a figure environment, an image-inclusion command or drawing code, so no image is delivered and the package declares no asset dependency.  Editorial formatting only, in addition: an AMS \texttt{amsart} preamble that restates the article's own declarations and macros used by the retained lines, namely the environments \texttt{lemma} on separate counters, together with the standard packages those lines need.  The retained environments print in this document as Lemma 1 in order of appearance, whereas the article prints the same items with the numbering of its own full document.  The complete native source of the article as distributed in the arXiv e-print for this exact version is bundled unchanged as \texttt{sources/210098/source0.tex}.
\begin{lemma}[Two braid identities]\label{lem:flype-identities}
For every integer $k\ge2$,
\begin{align}
\sigma_2
\bigl(
\sigma_1\sigma_2^k\sigma_1^{-k}\sigma_2^{-1}
\bigr)
\sigma_2^{-1}
&=
\sigma_1^{k-1}\sigma_2^{-(k-1)}\sigma_1\sigma_2^{-1},
\label{eq:flype-identity-one}\\
\sigma_1
\bigl(
\sigma_1^k\sigma_2^2\sigma_1^{-k-1}\sigma_2^{-1}
\bigr)
\sigma_1^{-1}
&=
\sigma_1^k\sigma_2^{-1}\sigma_1\sigma_2^{-k}.
\label{eq:flype-identity-two}
\end{align}
In particular, the braid inside the parentheses on the left-hand side of
each identity is conjugate to the braid on the corresponding right-hand side.
\end{lemma}

\begin{proof}
Put $\Delta=\sigma_1\sigma_2\sigma_1=\sigma_2\sigma_1\sigma_2$. The braid relation gives $\Delta\sigma_1^j=\sigma_2^j\Delta$ and $\Delta\sigma_2^j=\sigma_1^j\Delta$ for every $j\in\mathbb Z$. We shall also use $\sigma_2\sigma_1=\Delta\sigma_2^{-1}$, $\sigma_1^{-2}\Delta=\sigma_2\sigma_1\sigma_2^{-1}$, and $\sigma_1^{-1}\Delta\sigma_1^{-1}=\sigma_2$, which follow immediately from the two expressions for $\Delta$.

For the first identity, we compute
\begin{align*}
&\sigma_2
 \bigl(\sigma_1\sigma_2^k\sigma_1^{-k}\sigma_2^{-1}\bigr)
 \sigma_2^{-1}
=\sigma_2\sigma_1\sigma_2^k\sigma_1^{-k}\sigma_2^{-2}
=\Delta\sigma_2^{k-1}\sigma_1^{-k}\sigma_2^{-2}\\
&\qquad
=\sigma_1^{k-1}\Delta\sigma_1^{-k}\sigma_2^{-2}
=\sigma_1^{k-1}\sigma_2^{-k}\Delta\sigma_2^{-2}
=\sigma_1^{k-1}\sigma_2^{-k}\sigma_1^{-2}\Delta\\
&\qquad
=\sigma_1^{k-1}\sigma_2^{-k}
 \bigl(\sigma_2\sigma_1\sigma_2^{-1}\bigr)
=\sigma_1^{k-1}\sigma_2^{-(k-1)}\sigma_1\sigma_2^{-1}.
\end{align*}

For the second identity, we similarly obtain
\begin{align*}
&\sigma_1
 \bigl(\sigma_1^k\sigma_2^2\sigma_1^{-k-1}\sigma_2^{-1}\bigr)
 \sigma_1^{-1}
=\sigma_1^{k+1}\sigma_2^2\sigma_1^{-k-1}\sigma_2^{-1}\sigma_1^{-1}
=\sigma_1^k\Delta\sigma_1^{-1}\sigma_2
 \sigma_1^{-k-1}\sigma_2^{-1}\sigma_1^{-1}\\
&\qquad
=\sigma_1^k\sigma_2^{-1}\Delta\sigma_2
 \sigma_1^{-k-1}\sigma_2^{-1}\sigma_1^{-1}
=\sigma_1^k\sigma_2^{-1}\sigma_1\Delta
 \sigma_1^{-k-1}\sigma_2^{-1}\sigma_1^{-1}
=\sigma_1^k\sigma_2^{-1}\sigma_1\sigma_2^{-k-1}
 \Delta\sigma_2^{-1}\sigma_1^{-1}\\
&\qquad
=\sigma_1^k\sigma_2^{-1}\sigma_1\sigma_2^{-k-1}
 \sigma_1^{-1}\Delta\sigma_1^{-1}
=\sigma_1^k\sigma_2^{-1}\sigma_1\sigma_2^{-k-1}\sigma_2
=\sigma_1^k\sigma_2^{-1}\sigma_1\sigma_2^{-k}.
\end{align*}
This proves both identities directly from the braid relation.
\end{proof}

\end{document}
